
\subsubsection{4.1 Research Questions}

\begin{itemize}
\item \textbf{H-E1:} Does the MNIST-CNN zoo contain non-trivial permutation orbits and lack batch normalization?
\item \textbf{H-M1:} Do flat MLP encoders at matched capacity exhibit high permutation sensitivity (score > 0.3)?
\item \textbf{H-M2:} Do NFN encoders at matched capacity achieve near-zero permutation sensitivity (score < 0.1)?
\item \textbf{H-M3:} Does the symmetry spectrum hold (ρ(flat) < ρ(Deep Sets) < ρ(NFN)) with Δρ ≥ 0.05 and bootstrap 95\% CI lower bound > 0?
\end{itemize}

\subsubsection{4.2 Datasets}

Two variants of the Schurholt MNIST-CNN model zoo are used.

\textbf{H-E1 (orbit existence): \texttt{dataset_mnist_seed.pt} — seed-only variant.}
Contains 976 final-epoch checkpoints (filtered at training_iteration = 50 from 50,860 total records) of a smaller Conv(8)-Conv(8)-Conv(8)-FC-FC architecture. This variant was used for H-E1 because it was available at that stage of execution. It is architecturally distinct from the hyp_rand variant used for encoder training.

\textbf{H-M1/M2/M3 (encoder training and performance): \texttt{dataset_mnist_hyp_rand.pt} — hyp_rand variant.}
Contains 2,249 checkpoints of a BN-free CNN on MNIST with Conv(32)-Conv(64)-FC(128)-FC(10) architecture. Split: 1,589 / 322 / 338 (train / val / test). All Spearman ρ and Δρ results are computed exclusively on the hyp_rand test split (n = 338). Weight vector dimension: 2,464.

The orbit existence result (H-E1) characterizes the seed-only zoo (Conv(8) architecture), while the encoders were trained and evaluated on the hyp_rand zoo (Conv(32)-Conv(64) architecture). This constitutes an architectural mismatch between the orbit existence measurement and the encoder experiments; see Limitation L5 in Section 6.2.

\subsubsection{4.3 Evaluation Metrics}

Spearman ρ on the held-out test split (n = 338); Δρ = ρ(NFN) − ρ(flat MLP) with paired bootstrap 95\% CI (1,000 resamples); permutation sensitivity score (500 stratified pairs). The gate threshold for H-M3 is Δρ ≥ 0.05 with CI lower bound > 0.

\subsubsection{4.4 Implementation Details}

All experiments ran on a single CUDA GPU. The NFN checkpoint from H-M2 (best validation loss, epoch 114) was reused in H-M3. The Deep Sets encoder was trained fresh for H-M3 (best at epoch 39, val_loss = 0.0203). The trained flat MLP checkpoint from H-M1 was not found at the expected path at the time of H-M3 evaluation; the flat MLP was therefore evaluated with random (untrained) weights.

